\section*{Chapitre \ref{ch:pl:databases-and-sql} --- Bases de données et SQL}

\begin{solution}{exo:pl:databases-and-sql:1}
Une table des positions est une seconde copie de ce que disent les transactions, et deux copies peuvent diverger~; une requête sur les transactions ne le peut pas. Quand la requête est trop lente,
une table des positions peut être conservée comme cache, mais elle se reconstruit à partir des transactions et ne se corrige jamais d'elle-même.
\end{solution}

\begin{solution}{exo:pl:databases-and-sql:2}
Parce qu'une transaction corrigée a plusieurs versions, une par \omterm{def:pl:time-series-databases-and-the-alternatives:bucket}{intervalle de temps} système~: l'identifiant de la transaction se répète, et le début du \omterm{def:pl:databases-and-sql:system}{temps système}
permet de distinguer les versions.
\end{solution}

\begin{solution}{exo:pl:databases-and-sql:3}
500\,048 versions pour 500\,008 transactions~: chacune des 30 corrections de quantité et des 10 annulations clôt une version et en ajoute une~; les 8 enregistrements
tardifs sont des transactions nouvelles, avec une seule version chacune.
\end{solution}

\begin{solution}{exo:pl:databases-and-sql:4}
Sur les six ordres de deux lectures et deux écritures qui conservent la lecture de chaque écrivain avant son écriture, quatre placent les deux lectures avant l'une ou l'autre écriture~; alors
les deux écrivains calculent à partir de la même valeur de départ, et la seconde écriture écrase la première. Les 636 sur 1\,000 du modèle sont proches de quatre sixièmes.
\end{solution}

\begin{solution}{exo:pl:databases-and-sql:5}
Le mardi à 18~h en temps de validité sélectionne chaque transaction de lundi et de mardi, soit deux jours sur cinq ou 40~\% des lignes~; les lire par l'index
coûte à peu près autant que parcourir la table.
\end{solution}

\begin{solution}{exo:pl:databases-and-sql:6}
Seul le rapport corrigé~: la transaction tardive est valide à partir du mardi, mais son \omterm{def:pl:databases-and-sql:system}{temps système} commence le mercredi, donc le mardi à 18~h la base
ne la contenait pas.
\end{solution}

\begin{solution}{exo:pl:databases-and-sql:7}
Joindre la table des transactions à elle-même sur l'identifiant de transaction~: d'un côté, restreindre aux versions en vigueur au \omterm{def:pl:databases-and-sql:system}{temps système} du mardi à 18~h, de l'autre
aux versions en vigueur maintenant, et garder les couples dont la quantité ou la fin du temps de validité diffèrent. Sur la semaine du modèle, elle renvoie 40 transactions~: les 30 quantités corrigées
et les 10 annulations (requête \texttt{CHANGED} dans \texttt{pl\_tradedb.py}).
\end{solution}

\begin{solution}{exo:pl:databases-and-sql:8}
Une table d'audit enregistre les modifications que l'application a choisi de journaliser, dans son propre format~; reconstituer un état passé revient à les rejouer à rebours
sans erreur, y compris les modifications faites hors de l'application (un script, une correction manuelle). Avec le \omterm{def:pl:databases-and-sql:system}{temps système}, les états passés sont les données
elles-mêmes, interrogées comme le présent, et rien ne peut les modifier.
\end{solution}

\begin{solution}{pb:pl:databases-and-sql:1}
\begin{enumerate}
\item Que chaque instrument, chaque compte et chaque version de transaction soient identifiés une seule fois, et qu'aucune transaction ne renvoie
à un compte ou à un instrument qui n'existe pas.
\item Les faits sur les instruments et les comptes (devise, noms)~: ils vivent une seule fois dans leurs propres tables.
\item Pour que chaque copie de la base applique les mêmes modifications une seule fois, dans l'ordre, et puisse indiquer quel schéma elle possède.
\item Un \omterm{def:pl:time-series-databases-and-the-alternatives:bucket}{intervalle de temps} de validité et un \omterm{def:pl:time-series-databases-and-the-alternatives:bucket}{intervalle de temps} système.
\item Corrigée~: le \omterm{def:pl:databases-and-sql:system}{temps système} de la version courante est clos et la version corrigée insérée. Annulée~: de même, avec un \omterm{def:pl:time-series-databases-and-the-alternatives:bucket}{intervalle de temps} de validité
vide.
\item T1 lit une donnée, T2 la met à jour, et T1 écrit une valeur calculée à partir de sa lecture antérieure et valide~: la mise à jour de T2 est perdue.
\item Parce que sa lecture et son écriture sont des transactions distinctes~: chacune est sérialisable, mais le couple ne forme pas une seule transaction.
\item 636 lectures puis écritures~; aucune à l'intérieur d'une transaction~; aucune en une seule instruction.
\item Elle prend le verrou d'écriture au départ, si bien que le second écrivain est refusé jusqu'à la validation du premier, puis réessaie avec la nouvelle valeur.
\item Lorsque la nouvelle valeur est une fonction de l'ancienne que SQL sait exprimer, comme un incrément.
\item Un parcours de toute la table des transactions pour les trois.
\item L'index (compte, instrument) ramène la requête ponctuelle de 13.7~ms à 0.01~ms et laisse les autres inchangées~; l'index sur le temps de validité
ne change presque rien.
\item DuckDB~: rapport 9.8~ms contre 150, requête analytique 2.8~ms contre 148, requête ponctuelle 1.0~ms contre 13.7 sans index et 0.01 avec.
\item Chaque insertion et chaque correction doivent le mettre à jour aussi, et il occupe de l'espace.
\item Une copie en colonnes (DuckDB ici), rafraîchie à partir du magasin transactionnel.
\item \textbf{Résultat nommé.} Le rapport des positions de mardi à 18~h est reproduit exactement à partir de la table bitemporelle (valide mardi à 18~h, système
mardi à 18~h) et diffère du rapport corrigé (système vendredi) sur 48 positions parmi 10\,000, une par correction~: 30 quantités, 10
annulations et 8 enregistrements tardifs. L'index (compte, instrument) accélère la requête ponctuelle environ mille fois (de 13.7~ms à 0.01~ms)~; l'index sur le
temps de validité ne gagne presque rien sur aucune requête.
\item Le \omterm{def:pl:databases-and-sql:system}{temps système} est estampillé par la base lorsqu'elle enregistre un fait et ne peut pas être antidaté~; le temps de connaissance est le moment où le fait est devenu disponible,
fourni avec la donnée.
\item Ce que la base croyait avant chaque correction, et donc chaque rapport produit avant elle.
\item Qu'elle peut répondre pour n'importe quel \omterm{def:pl:databases-and-sql:system}{temps système} passé.
\item Ce qu'elle croyait, et quand.
\end{enumerate}
\end{solution}

\begin{solution}{iq:pl:databases-and-sql:1}
Un groupe de lectures et d'écritures appliqué comme une unité~: atomique (tout ou rien), cohérent (d'un état valide à un état valide), isolé (aucun travail partiel
n'est vu des autres) et durable (une transaction validée survit à une panne).
\iqlookfor{Les quatre, et ce que chacune protège.}
\end{solution}

\begin{solution}{iq:pl:databases-and-sql:2}
Les deux lisent l'ancienne valeur et écrivent leur propre somme, et une écriture écrase l'autre. Envoyer les incréments en une seule instruction, ou faire la lecture et
l'écriture dans une même transaction qui prend le verrou d'écriture au départ, avec une nouvelle tentative.
\iqlookfor{Lecture-modification-écriture en une seule transaction.}
\end{solution}

\begin{solution}{iq:pl:databases-and-sql:3}
Lire le plan~: un parcours complet d'une grande table pour une requête qui sélectionne peu de lignes réclame un index sur les colonnes filtrées~; une requête qui lit
la plupart des lignes n'en réclame pas. Mesurer avant et après, et compter le coût sur les écritures.
\iqlookfor{Sélectivité, plan, mesure.}
\end{solution}

\begin{solution}{iq:pl:databases-and-sql:4}
Conserver le \omterm{def:pl:databases-and-sql:system}{temps système}~: les corrections closent des versions et en insèrent de nouvelles, sans jamais mettre à jour~; produire les rapports à partir de requêtes au \omterm{def:pl:databases-and-sql:system}{temps système} voulu, et enregistrer
le \omterm{def:pl:databases-and-sql:system}{temps système} utilisé par chaque rapport.
\iqlookfor{Tables bitemporelles, jamais de mise à jour sur place.}
\end{solution}

\begin{solution}{iq:pl:databases-and-sql:5}
Un magasin transactionnel avec clés, migrations et transactions bitemporelles pour l'enregistrement~; les positions sous forme de requêtes ou de caches reconstruits~; une copie en colonnes pour
le reporting et la recherche, rafraîchie à partir de celui-ci~; des index choisis à partir des plans~; chaque rapport estampillé de son \omterm{def:pl:databases-and-sql:system}{temps système}.
\iqlookfor{Magasin de référence OLTP, copie OLAP, time travel.}
\end{solution}
